\documentclass[../DoAn.tex]{subfiles}
\begin{document}

\begin{center}
    \Large{\textbf{ABSTRACT}}\\
\end{center}
\vspace{1cm}

Multi-Agent Pathfinding (MAPF) requires simultaneous path planning for multiple agents
navigating from their start positions to designated goals without collisions. In real-time
strategy games, this problem becomes particularly challenging due to continuously changing
environments, demanding agents to replan within milliseconds each frame. Traditional
single-agent solutions such as A* do not inherently prevent inter-agent collisions, while
globally optimal approaches like CBS impose computational costs too high for real-time
environments.

This thesis studies and implements three MAPF algorithm tiers within a 2D multiplayer
tank game built on Unity Engine: (1) single-agent A* as a baseline, (2) PIBT (Priority
Inheritance with Backtracking) implemented directly in Unity~C\# for collision-free
multi-agent coordination, and (3) PIBT C++/TCP integrating an external planning server
over TCP, enabling the use of high-performance C++ algorithms alongside Unity's physics.
The game uses standard MovingAI MAPF benchmark maps, ensuring results are comparable
with community research.

Key contributions include: (i) a runtime map loading system that constructs grid and
physics colliders from the \texttt{.map} format without baking into the Unity scene;
(ii) A* and PIBT implementations in C\# fully integrated with Unity physics and the
shared tank control pipeline; (iii) a TCP protocol bridging Unity and an external
PIBT C++ server; and (iv) an automated backtest system supporting static and
dynamic-obstacle environments, exporting per-run CSV metrics. Experiments across
5 maps, 3 algorithms, and 6 repetitions per combination (90 runs per condition)
show all three algorithms successfully destroy the Eagle target on 4 of 5 maps;
PIBT exhibits higher replanning counts than A* due to its cooperative coordination
mechanism; and enabling dynamic obstacles increases replanning by 30--50\% on
average while the overall mission success rate remains stable.

\end{document}
